\chapter{Stage A Molnar Single-Notch Benchmark}

\section{Objective and Decision Boundary}

Stage A establishes whether the Molnar and Gravouil single-notch phase-field
benchmark is a defensible baseline for later Abaqus adaptive-remeshing work.
The aim is deliberately narrower than the full thesis workflow: reproduce an
established staggered phase-field implementation, quantify the reconstructed
single-notch response, and decide whether the evidence is strong enough to
authorize Stage B uniform-reference studies.

The current Stage A status is:

\begin{quote}
\texttt{candidate v2: technically passed}\\
\texttt{scientific review: incomplete}\\
\texttt{SDV15: completed-increment irreversibility violation}\\
\texttt{Gate A3: reference\_data\_insufficient}\\
\texttt{new solution run: not authorized}
\end{quote}

This chapter therefore treats candidate v2 as a technically reproducible and
useful reproduction baseline, but not as a Gate A3 pass. Supervisor approval is
required before any Stage B, uniform-reference, MISESERI, remeshing, or state
transfer work begins.

\section{Benchmark Choice}

The Molnar single-notch benchmark was selected because it connects directly to
the thesis requirement for an Abaqus UEL/UMAT phase-field workflow. The source
package contains a staggered displacement/phase implementation, a visualization
layer through state variables, and single-notch examples suitable for checking
reaction force, crack-path geometry, and phase-field/history behavior before
introducing adaptive remeshing.

The benchmark is also a useful stress test for thesis discipline. It separates
technical execution from scientific validation: a model can compile, solve, and
produce a plausible crack while still failing a strict scientific gate because
the RF--U reference, tolerances, or irreversibility interpretation remain
unresolved.

\section{Formulation Convention}

The phase-field convention used throughout Stage A is:

\begin{itemize}
\item \(d=0\): intact material.
\item \(d=1\): fully broken material.
\end{itemize}

Candidate v2 preserves the published Molnar source formulation. The generated
source copy \path{models/generated/molnar_gravouil_2017/paper_matched_single_notch_v2/SingleNotch_v2.for}
is copied from the preserved source with only the physical-element count
\texttt{N\_ELEM=33852} updated for the reconstructed mesh. The original source
is preserved at
\path{models/baseline_original/molnar_gravouil_2017/02_Single_Notch_Tension/SingleNotch.for}.

\section{Staggered Phase/Displacement Method}

The implementation follows a staggered solution structure with displacement
and phase-field layers. The displacement field is solved in the mechanical
layer, and the phase-field layer updates the scalar damage variable using the
stored history information. Visualization variables are copied into a companion
UMAT/CPS4 layer; those visualization fields are evidence outputs, not separate
governing fields.

For thesis use, three distinctions must remain visible:

\begin{itemize}
\item The published formulation is the Molnar and Gravouil staggered
phase-field UEL/UMAT implementation.
\item The reconstructed implementation is the project candidate-v2
paper-matched deck generated from documented geometry, loading, material, and
mesh assumptions.
\item The approximate Fig.~7 reference is a local digitization of the published
curve, not exact author data and not a supervisor-approved final tolerance
basis.
\end{itemize}

\section{Geometry, Material, and Loading}

Candidate v2 reconstructs a \(1.0~\mathrm{mm}\times1.0~\mathrm{mm}\)
single-edge-notched plate with a left-edge split-node notch of length
\(0.5~\mathrm{mm}\) at \(y=0\). The model uses plane strain and unit
thickness. Material parameters are \(E=210~\mathrm{kN/mm^2}\), \(\nu=0.3\),
\(G_c=0.0027~\mathrm{kN/mm}\), and \(l_c=0.0075~\mathrm{mm}\).

The load angle is \(90^\circ\). The loading is applied in two direct-increment
steps:

\begin{center}
\begin{tabular}{lrrr}
\toprule
Step & Start \(U_2\) (mm) & End \(U_2\) (mm) & Increment (mm) \\
\midrule
Step 1 & 0.0000 & 0.0050 & \(1.0\times10^{-5}\) \\
Step 2 & 0.0050 & 0.0067 & \(1.0\times10^{-5}\) \\
\bottomrule
\end{tabular}
\end{center}

\section{Paper-Matched Reconstruction}

The reconstruction is documented in the project configuration and derived
reference records. Candidate v2 restores the intended loading arithmetic,
implements a split-node notch, preserves the layered UEL/UMAT architecture,
and uses the source property layout from the preserved Molnar files. Static
validation passed, and the mesh-quality preflight classified the model as
\texttt{mesh\_quality\_preflight\_pass}.

The mesh has \(33852\) physical elements and \(101556\) layered elements. The
local mesh size is \(h=0.001~\mathrm{mm}\), giving \(h/l=0.1333333333\). The
maximum element size is \(0.025~\mathrm{mm}\), the maximum neighboring-size
ratio is \(1.5\), and the maximum aspect ratio is \(25.000000000001368\). The
high-aspect-ratio elements are documented as reconstruction limitations outside
the fracture-process corridor.

\section{Technical Validation}

The single authorized serial HPC baseline job was
\texttt{1374864.mmaster02}, submitted from revision
\texttt{711dd495bdcb830d695f9d7e56283316c9d417d5}. It completed with PBS
\texttt{Exit\_status=0}; Abaqus returned code zero; the \texttt{.sta},
\texttt{.msg}, \texttt{.dat}, and scratch \texttt{.odb} outputs existed; and
the \texttt{.sta} reported successful analysis completion. The technical
classification is \texttt{paper\_matched\_v2\_technical\_pass}.

The run used one CPU, \(32~\mathrm{GB}\) requested memory, and a 24-hour
walltime request. Recorded resource use was walltime \(00{:}38{:}38\), CPU
time \(00{:}35{:}52\), CPU percent \(95\), memory \(2970760~\mathrm{kb}\),
virtual memory \(3565868~\mathrm{kb}\), and an Abaqus-reported peak memory of
about \(2~\mathrm{GB}\).

\section{Numerical Results}

\begin{longtable}{p{0.40\linewidth}p{0.25\linewidth}p{0.27\linewidth}}
\caption{Stage A candidate-v2 numerical evidence. Values are drawn from
committed extraction, scientific-review, and targeted-diagnostic evidence.}\\
\toprule
Metric & Value & Evidence \\
\midrule
\endfirsthead
\toprule
Metric & Value & Evidence \\
\midrule
\endhead
Peak RF2 & \(0.761702~\mathrm{kN}\) & \path{runs/hpc/paper_matched_single_notch_v2/scientific_review/fig7_comparison_metrics.json} \\
Peak displacement & \(U_2=0.006110~\mathrm{mm}\) & \path{runs/hpc/paper_matched_single_notch_v2/scientific_review/fig7_comparison_metrics.json} \\
Peak-force error & \(0.064519\) relative & \path{runs/hpc/paper_matched_single_notch_v2/scientific_review/FIG7_COMPARISON_AUDIT.md} \\
Peak-displacement error & \(0.041257\) relative & \path{runs/hpc/paper_matched_single_notch_v2/scientific_review/FIG7_COMPARISON_AUDIT.md} \\
Pre-peak RMSE & \(0.044136~\mathrm{kN}\) & \path{runs/hpc/paper_matched_single_notch_v2/scientific_review/FIG7_COMPARISON_AUDIT.md} \\
Post-peak RMSE & \(0.348093~\mathrm{kN}\) & \path{runs/hpc/paper_matched_single_notch_v2/scientific_review/FIG7_COMPARISON_AUDIT.md} \\
Full-overlap NRMSE & \(0.245705\) & \path{runs/hpc/paper_matched_single_notch_v2/scientific_review/fig7_comparison_metrics.json} \\
Tangent-stiffness error & about \(+51.4\%\) & \path{runs/hpc/paper_matched_single_notch_v2/scientific_review/SCIENTIFIC_DECISION.md} \\
Common-domain area error & \(+0.193324\) relative & \path{runs/hpc/paper_matched_single_notch_v2/scientific_review/fig7_comparison_metrics.json} \\
Crack extension, SDV15 \(\ge0.80\) & \(0.0555~\mathrm{mm}\) beyond notch tip; \(0.0570~\mathrm{mm}\) total connected path & \path{runs/hpc/paper_matched_single_notch_v2/scientific_review/SCIENTIFIC_DECISION.md} \\
Crack extension, SDV15 \(\ge0.90\) & \(0.0525~\mathrm{mm}\) beyond notch tip; \(0.0530~\mathrm{mm}\) total connected path & \path{runs/hpc/paper_matched_single_notch_v2/scientific_review/SCIENTIFIC_DECISION.md} \\
Crack extension, SDV15 \(\ge0.95\) & \(0.0505~\mathrm{mm}\) beyond notch tip; \(0.0510~\mathrm{mm}\) total connected path & \path{runs/hpc/paper_matched_single_notch_v2/scientific_review/SCIENTIFIC_DECISION.md} \\
Crack extension, SDV15 \(\ge0.99\) & \(0.0465~\mathrm{mm}\) beyond notch tip; \(0.0460~\mathrm{mm}\) total connected path & \path{runs/hpc/paper_matched_single_notch_v2/scientific_review/SCIENTIFIC_DECISION.md} \\
SDV15 maximum overshoot & \(1.005600\) & \path{runs/hpc/paper_matched_single_notch_v2/scientific_review/sdv_irreversibility_metrics.json} \\
Accepted-increment decreasing transitions & \(2184\) & \path{runs/hpc/paper_matched_single_notch_v2_sdv15_diagnostic_r2/RUN_SUMMARY.md} \\
Transitions above \(10^{-6}\) & \(2131\) & \path{runs/hpc/paper_matched_single_notch_v2_sdv15_diagnostic_r2/evidence/1375028.mmaster02/postprocessing_completed_increment_replay_time_aligned/severity_audit/sdv15_completed_increment_severity_summary.json} \\
Transitions above \(10^{-5}\) & \(1362\) & \path{runs/hpc/paper_matched_single_notch_v2_sdv15_diagnostic_r2/evidence/1375028.mmaster02/postprocessing_completed_increment_replay_time_aligned/severity_audit/sdv15_completed_increment_severity_summary.json} \\
Maximum decrease & \(2.2384088238425193\times10^{-4}\) & \path{runs/hpc/paper_matched_single_notch_v2_sdv15_diagnostic_r2/RUN_SUMMARY.md} \\
Mean decrease & \(2.8434597987446906\times10^{-5}\) & \path{runs/hpc/paper_matched_single_notch_v2_sdv15_diagnostic_r2/RUN_SUMMARY.md} \\
Median decrease & \(1.3940791172561973\times10^{-5}\) & \path{runs/hpc/paper_matched_single_notch_v2_sdv15_diagnostic_r2/RUN_SUMMARY.md} \\
Unique affected element/IP locations & \(138\) & \path{runs/hpc/paper_matched_single_notch_v2_sdv15_diagnostic_r2/RUN_SUMMARY.md} \\
SDV16 decreases & \(0\) & \path{runs/hpc/paper_matched_single_notch_v2_sdv15_diagnostic_r2/RUN_SUMMARY.md} \\
Diagnostic RF--U perturbation & maximum normalized RF2 difference \(6.54442468760855\times10^{-13}\); maximum U2 difference \(4.992967844730245\times10^{-15}~\mathrm{mm}\) & \path{runs/hpc/paper_matched_single_notch_v2_sdv15_diagnostic_r2/RUN_SUMMARY.md} \\
Solver walltime and memory & baseline PBS walltime \(00{:}38{:}38\), memory \(2970760~\mathrm{kb}\), Abaqus peak memory about \(2~\mathrm{GB}\) & \path{runs/hpc/paper_matched_single_notch_v2/scientific_review/solver_resource_metrics.json} \\
\bottomrule
\end{longtable}

\section{RF--U Interpretation}

The RF--U comparison is mixed. The peak force scale is close enough to justify
continued review, but the full curve is not within a strict provisional
working gate. The pre-reference-peak RMSE is \(0.044136~\mathrm{kN}\), while
the post-reference-peak RMSE is \(0.348093~\mathrm{kN}\). This makes the
post-peak mismatch the dominant scientific limitation. The comparison uses the
approximate digitized Fig.~7 red dashed \(l_c=0.0075~\mathrm{mm}\) reference,
which has estimated digitization uncertainty of about
\(\pm0.00004~\mathrm{mm}\) and \(\pm0.015~\mathrm{kN}\).

\section{Crack-Path Interpretation}

The final high-damage band is connected and approximately horizontal at
thresholds from \(\mathrm{SDV15}\ge0.80\) through
\(\mathrm{SDV15}\ge0.99\). At the provisional
\(\mathrm{SDV15}\ge0.95\) threshold, the final matched contour contains
193 element-mean damaged elements and a connected high-damage extension of
about \(0.0505~\mathrm{mm}\) beyond the notch tip. This supports qualitative
crack-direction agreement, but it remains threshold-dependent and does not
close Gate A3 by itself.

\section{SDV15, SDV16, and Irreversibility}

In the project convention, SDV15 is the phase-field visualization variable and
SDV16 is the history-like variable used to monitor monotonic driving behavior.
The original ODB scan found SDV16 monotone and SDV15 nonmonotone in retained
visualization frames. A later targeted diagnostic r2 run was non-intrusive by
RF--U comparison and enabled a no-solution completed-increment replay of the
existing trace rows.

The completed-increment replay kept the last U1 stage-101 call for each
\((\mathrm{KSTEP}, \mathrm{KINC}, \mathrm{physical~element},
\mathrm{source-storage~IP})\). It processed 627304 trace rows, retained
209152 U1 stage-101 rows, formed 101840 final increment states over 152
element/IP sequences, and found 2184 unique completed-increment SDV15
decreasing transitions. The severity audit classifies this as
\texttt{sdv15\_completed\_increment\_irreversibility\_violation} under the
provisional \(10^{-6}\) materiality tolerance. SDV16 has zero decreases over
the same final-increment sequences.

\section{Scientific Limitations}

The current evidence supports the following restrained interpretation:

\begin{itemize}
\item Candidate v2 is technically reproducible.
\item Crack propagation direction is qualitatively correct.
\item RF--U agreement is mixed, and post-peak disagreement is substantial.
\item SDV16 remains monotone in the checked sequences.
\item SDV15 decreases between accepted increments.
\item The targeted diagnostic instrumentation is non-intrusive.
\item Strict pointwise phase-field irreversibility is not satisfied.
\item Candidate v2 remains a useful reproduction baseline for supervisor
review.
\item Supervisor approval is required before Stage B or any downstream
adaptive-remeshing work.
\end{itemize}


\section{Supplementary lc = 0.015 mm mesh h-convergence (RF--U)}
\label{sec:lc015-hconv-thesis}

For the exact author-supplied supplementary single-notch model (\(l_c=0.015\,\mathrm{mm}\)), three serial meshes were solved and postprocessed for RF--U only.
Successive peak-force changes are approximately \(4.0\%\) (H0\(\to\)H1) and \(0.47\%\) (H1\(\to\)H2-PUB); pre-peak curve NRMSE between H1 and H2-PUB is about \(0.11\%\), while post-peak residual remains larger (\(\sim 20\%\) NRMSE).
After supervisor Decisions~\textbf{1A} and \textbf{2B}, mesh roles are frozen as:
H2-PUB for fine RF--U validation only; H1 for production/thesis/report simulations and the Stage~C comparison reference; H0 for fast testing and the first MISESERI implementation trials.
Crack-path/contour mesh convergence remains deferred (Decision~\textbf{2B}) and does not block Stage~C preparation.
Evidence: \path{runs/hpc/molnar_lc015_h_convergence/comparison/H_CONVERGENCE_SCIENTIFIC_REVIEW.md};
\path{docs/decisions/MOLNAR_GATE_A3_SUPERVISOR_DECISION_1A_2B.md};
\path{docs/decisions/MESH_USE_POLICY.md}.

\section{Gate A3 Boundary}

Gate A3 is \textbf{conditionally accepted for RF--U validation use} after Decisions~\textbf{1A} and \textbf{2B}
(\texttt{gate\_a3\_conditionally\_accepted\_rf\_u}; \texttt{contour\_validation\_deferred}).
Full unconditional closure of every historical Stage~A claim is not asserted; residual items may still retain a
\texttt{reference\_data\_insufficient} narrative where appropriate.
Candidate v2 has not been reclassified as a full scientific pass by these decisions.
Stage~C MISESERI \textbf{preparation} is authorized; HPC submission, MISESERI execution, multicore qualification,
and state-transfer work remain unauthorized without explicit further approval.
Stage~C preparation plan: \path{docs/studies/STAGE_C_MISESERI_PREPARATION_PLAN.md}.
